\chapter{Теоретические сведения}
\label{ch:theory}

\section{HTML как язык разметки}

\textbf{HTML} — декларативный язык разметки гипертекста, разработанный
Тимом Бернерсом-Ли в 1991~году и поддерживаемый консорциумом W3C
(текущая редакция стандарта --- HTML~Living~Standard, WHATWG). Документ
состоит из вложенных \emph{элементов}; каждый элемент описывается тегом
вида \texttt{<имя~атрибуты>...</имя>} (парный тег) или
\texttt{<имя~атрибуты>} (одиночный, например \texttt{<img>},
\texttt{<br>}, \texttt{<hr>}).

\textbf{Атрибут} — пара <<имя=\,значение>> внутри открывающего тега,
уточняющая поведение элемента: \texttt{href} у ссылки,
\texttt{src}/\texttt{alt} у изображения, \texttt{class}/\texttt{id} у
любого элемента, \texttt{colspan}/\texttt{rowspan} у ячейки таблицы.

Минимальная структура HTML-документа:
\begin{lstlisting}[language=HTML]
<!DOCTYPE html>
<html lang="ru">
<head>
  <meta charset="UTF-8">
  <title>Заголовок страницы</title>
</head>
<body>
  <!-- содержимое -->
</body>
</html>
\end{lstlisting}

\texttt{<!DOCTYPE~html>} включает стандартный (не quirks) режим
рендеринга, \texttt{<html~lang>} задаёт язык документа (важно для
проверки орфографии и поисковой оптимизации),
\texttt{<meta~charset="UTF-8">} объявляет кодировку.

\section{Группировка контента и текстовая разметка}

Семантические контейнеры группируют логически связанные блоки
страницы:
\begin{itemize}
  \item \texttt{<header>} --- шапка раздела или документа;
  \item \texttt{<section>} --- смысловой раздел с собственной темой;
  \item \texttt{<article>} --- самостоятельная публикация (статья,
        пост);
  \item \texttt{<aside>} --- врезка, не связанная с основным
        потоком;
  \item \texttt{<footer>} --- подвал раздела или документа.
\end{itemize}

Текстовая разметка:
\begin{itemize}
  \item \texttt{<h1>}--\texttt{<h6>} --- заголовки шести уровней
        (от самого крупного к самому мелкому); единственный
        \texttt{<h1>} на странице обычно используется для главного
        заголовка документа;
  \item \texttt{<p>} --- абзац;
  \item \texttt{<blockquote>} --- блочная цитата, типично включает
        внутрь себя один или несколько \texttt{<p>};
  \item \texttt{<strong>}, \texttt{<em>} --- логическое выделение
        (важность и интонационный акцент соответственно);
  \item \texttt{<span>} --- безразмерный встраиваемый контейнер
        без собственной семантики, нужен для стилизации
        фрагмента текста через \texttt{class};
  \item \texttt{<a~href="URL">текст</a>} --- гиперссылка;
        атрибут \texttt{target="\_blank"} открывает её в новой
        вкладке;
  \item \texttt{<img~src="URL"~alt="описание">} --- одиночный тег
        изображения; атрибут \texttt{alt} обязателен для
        доступности (озвучивается скринридером, показывается
        вместо картинки при ошибке загрузки).
\end{itemize}

\section{Списки}

В HTML три вида списков:
\begin{itemize}
  \item \texttt{<ul>} --- неупорядоченный (маркированный) список;
        элементы внутри --- \texttt{<li>}; маркер по умолчанию ---
        закрашенный кружок;
  \item \texttt{<ol>} --- упорядоченный (нумерованный) список;
        элементы --- \texttt{<li>}; маркер по умолчанию ---
        арабская цифра с точкой; атрибуты \texttt{start},
        \texttt{type}, \texttt{reversed} управляют нумерацией;
  \item \texttt{<dl>} --- описательный список из пар
        <<термин -- определение>>: \texttt{<dt>} для термина,
        \texttt{<dd>} для определения.
\end{itemize}

Списки могут быть вложенными --- внутри \texttt{<li>} разрешается
размещать ещё один \texttt{<ul>} или \texttt{<ol>}; такая структура
используется, например, в оглавлениях с подразделами.

\section{Таблицы}

Таблица в HTML строится из иерархии тегов:
\begin{itemize}
  \item \texttt{<table>} --- корневой контейнер;
  \item \texttt{<thead>}, \texttt{<tbody>}, \texttt{<tfoot>} ---
        семантические группы строк (шапка, тело, подвал);
  \item \texttt{<tr>} --- строка таблицы;
  \item \texttt{<th>} --- ячейка-заголовок (по умолчанию жирный
        текст и выравнивание по центру);
  \item \texttt{<td>} --- обычная ячейка с данными.
\end{itemize}

Объединение ячеек выполняется атрибутами:
\begin{itemize}
  \item \texttt{colspan="N"} --- ячейка занимает N колонок по
        горизонтали;
  \item \texttt{rowspan="N"} --- ячейка занимает N строк по
        вертикали.
\end{itemize}

Если в строке часть ячеек объединена через \texttt{colspan}, то
суммарное число колонок (с учётом объединений) во всех строках
должно совпадать --- иначе таблица <<разваливается>>: появляются
лишние колонки в одних строках и не хватает ячеек в других. Это
типичная ошибка вёрстки, которая встречается в задании~4 блока~1.2
текущей лабораторной работы.

\section{CSS: базовая модель}

\textbf{CSS} --- декларативный язык, описывающий правила оформления
HTML-элементов. Правило состоит из \emph{селектора} и блока
\emph{объявлений}:
\begin{lstlisting}[language=CSS]
селектор {
  свойство: значение;
  свойство: значение;
}
\end{lstlisting}

Способы подключения CSS к документу:
\begin{itemize}
  \item \textbf{inline} --- через атрибут \texttt{style="..."}
        конкретного элемента; используется редко (плохо
        масштабируется);
  \item \textbf{внутренний} --- блок \texttt{<style>~...~</style>}
        в \texttt{<head>}; подходит для одиночных страниц;
  \item \textbf{внешний} --- через
        \texttt{<link~rel="stylesheet"~href="...">}; основной способ
        для проектов из нескольких страниц --- стили выносятся в
        отдельный файл и переиспользуются.
\end{itemize}

В этой работе используется именно внешний способ: в каталоге каждого
задания --- собственный \texttt{styles.css}, подключаемый из
\texttt{index.html}.

\section{CSS-селекторы}

Базовые селекторы:
\begin{itemize}
  \item \textbf{по тегу} --- \texttt{p}, \texttt{strong} ---
        выбирает все элементы соответствующего типа;
  \item \textbf{по классу} --- \texttt{.main-header} ---
        выбирает все элементы с \texttt{class="main-header"};
  \item \textbf{по идентификатору} --- \texttt{\#footer} ---
        выбирает единственный элемент с \texttt{id="footer"};
  \item \textbf{универсальный} --- \texttt{*} --- любой элемент.
\end{itemize}

Комбинаторы:
\begin{itemize}
  \item \texttt{A~B} --- \emph{потомок}: B на любой глубине
        внутри A (например, \texttt{blockquote~span} --- любой
        \texttt{<span>} внутри любого \texttt{<blockquote>});
  \item \texttt{A~>~B} --- \emph{прямой потомок}: B
        непосредственно внутри A, без промежуточных
        контейнеров;
  \item \texttt{A~+~B} --- следующий за A элемент B;
  \item \texttt{A~\textasciitilde~B} --- все последующие соседи B
        после A.
\end{itemize}

Составные селекторы пишутся без пробелов: \texttt{p.lead} ---
абзацы с классом \texttt{lead}; \texttt{a:hover} --- ссылка в
состоянии наведения курсора.

При конфликте правил (несколько селекторов задают одно свойство
одному элементу) побеждает более \emph{специфичный}: ID-селекторы
сильнее классов, классы сильнее тегов. При равной специфичности
выигрывает правило, объявленное позднее в порядке загрузки таблиц
стилей.

\section{Типографические свойства CSS}

Свойства, использованные в задачах текущей лабораторной работы:
\begin{itemize}
  \item \texttt{color: <цвет>} --- цвет текста (именованный,
        HEX, RGB, HSL);
  \item \texttt{font-size: <размер>} --- кегль шрифта
        (\texttt{px}, \texttt{em}, \texttt{rem}, \texttt{\%});
  \item \texttt{font-family: <семейство>} --- семейство шрифта
        (\texttt{serif} --- с засечками, \texttt{sans-serif} ---
        без засечек, \texttt{monospace} --- моноширинный,
        либо конкретное имя гарнитуры);
  \item \texttt{font-style: italic | normal} --- курсивное /
        прямое начертание;
  \item \texttt{font-weight: bold | normal | 100..900} ---
        жирность;
  \item \texttt{text-decoration: underline | line-through | none} ---
        декоративные линии вокруг текста;
  \item \texttt{list-style-type: disc | circle | square |
        decimal | none} --- тип маркера списка;
  \item \texttt{background-image: url("...")} --- фоновое
        изображение элемента;
  \item \texttt{background-color: <цвет>} --- фоновая заливка.
\end{itemize}

\endinput
